\subsection{AI Literacy Outside the Classroom}
\label{sec:8.2.3}
Most AI literacy work assumes someone who has already decided to learn: a course to enroll in, a curriculum to sit through, a public MOOC to find. That reaches the people who came looking. The harder and more useful problem is everyone else — reaching people where they already are, in a policy debate, a town hall, a local news story, without their having gone out after it.

The engagement that does most work is the kind attached to a live question. The Ada Lovelace Institute's Citizens' Biometrics Council brought fifty members of the UK public together to deliberate on how biometric technology should be regulated, alongside the first national survey of UK public opinion on the technology \autocite{adalovelaceinstitute2021biometrics}. The specificity is what made it work. Participants were not asked what they thought about AI, a question almost nobody can answer usefully, but about facial recognition — a thing with a location, an operator, and a consequence.

That is the lesson for public literacy generally, and it cuts against most of what gets funded under the heading. General AI awareness produces general opinions, which are worth little to a regulator and nothing to a court. A public able to govern a technology is a public that has been asked about a particular use of it.
